% ============================================================
%  Chapter 2 -- Metric Spaces, 2.15-2.30
% ============================================================
\rsection{Metric Spaces}

\begin{signpost}[Why this section is here]
Chapter 1 ended by noting that $\R^k$ has a distance function satisfying three
properties (1.37(a),(b),(f)). This section takes those three properties as the
\emph{definition} of a space and develops everything that follows from them
alone. The gain is reach: every theorem proved here holds for $\R$, for $\C$,
for $\R^k$, for subsets of any of these, and later for spaces of functions
(7.14, 11.35).

\smallskip
Definition 2.18 is the hinge of the chapter --- ten pieces of vocabulary in one
block, arriving with no motivation. Do not try to absorb them all. The two that
carry the weight are \textbf{limit point} (b) and \textbf{open} (f); ``closed''
is defined from the first, and everything in 2.19--2.30 is bookkeeping about
how the two interact. The rest of the list is there so Rudin need not
interrupt himself later.

\smallskip
Read 2.29--2.30 with attention even though they look technical. They say that
``open'' is not a property a set has by itself but one it has \emph{relative to
an ambient space}, and that ambiguity causes more confusion in Chapter 4 than
any other point in the book.
\end{signpost}

\ritem{2.15}{Definition}
A set $X$, whose elements we shall call \emph{points}, is said to be a
\emph{metric space} if with any two points $p$ and $q$ of $X$ there is
associated a real number $d(p,q)$, called the \emph{distance} from $p$ to $q$,
such that
\begin{enumerate}[label=(\alph*), leftmargin=2.2em, itemsep=1pt, topsep=3pt]
  \item $d(p,q) > 0$ if $p \ne q$; $d(p,p) = 0$;
  \item $d(p,q) = d(q,p)$;
  \item $d(p,q) \le d(p,r) + d(r,q)$ for any $r \in X$.
\end{enumerate}
Any function with these three properties is called a \emph{distance function},
or a \emph{metric}.\kn{Compare 1.37: (a) here is 1.37(a),(b); (b) follows from
1.37(c) with $\alpha = -1$; (c) is 1.37(f). Rudin verified all three for $\R^k$
one chapter ago, which is why 2.16 can assert the example in a single line.
Note what is \emph{absent}: no addition, no scalar multiplication, no inner
product. A metric space has distance and nothing else.}

\ritem{2.16}{Examples}
The most important examples of metric spaces, from our standpoint, are the
euclidean spaces $\R^k$, especially $\R^1$ (the real line) and $\R^2$ (the
complex plane); the distance in $\R^k$ is defined by
\begin{equation}
  d(\mathbf{x},\mathbf{y}) = |\mathbf{x}-\mathbf{y}|
  \qquad (\mathbf{x},\mathbf{y} \in \R^k). \tag{2.19}
\end{equation}
By Theorem 1.37, the conditions of Definition 2.15 are satisfied by (2.19).

It is important to observe that every subset $Y$ of a metric space $X$ is a
metric space in its own right, with the same distance function. For it is clear
that if conditions (a) to (c) of Definition 2.15 hold for $p,q,r \in X$, they
also hold if we restrict $p,q,r$ to lie in $Y$.\kn{This one-sentence
observation is what makes 2.29--2.30 necessary. Because \emph{every} subset is
itself a metric space, the words ``open'' and ``closed'' become ambiguous the
moment more than one ambient space is in play --- and from Chapter 4 onward
there always is.}

Thus, every subset of a euclidean space is a metric space. Other examples are
the spaces $\mathscr{C}(K)$ and $\mathscr{L}^2(\mu)$, which are discussed in
Chaps.\ 7 and 11, respectively.

\ritem{2.17}{Definition}
By the \emph{segment} $(a,b)$ we mean the set of all real numbers $x$ such that
$a < x < b$.

By the \emph{interval} $[a,b]$ we mean the set of all real numbers $x$ such
that $a \le x \le b$.

Occasionally we shall also encounter ``half-open intervals'' $[a,b)$ and
$(a,b]$; the first consists of all $x$ such that $a \le x < b$, the second of
all $x$ such that $a < x \le b$.

If $a_i < b_i$ for $i = 1,\dots,k$, the set of all points
$\mathbf{x} = (x_1,\dots,x_k)$ in $\R^k$ whose coordinates satisfy the
inequalities $a_i \le x_i \le b_i$ $(1 \le i \le k)$ is called a
\emph{$k$-cell}. Thus a $1$-cell is an interval, a $2$-cell is a rectangle,
etc.\kw{Rudin's ``segment'' means open interval and his ``interval'' means
closed. Almost no other book uses these words this way, and he uses them
constantly. A $k$-cell is always \emph{closed} --- the inequalities are
$\le$ --- which is why 2.40 can prove they are compact.}

If $\mathbf{x} \in \R^k$ and $r > 0$, the \emph{open} (or \emph{closed})
\emph{ball} $B$ with center at $\mathbf{x}$ and radius $r$ is defined to be the
set of all $\mathbf{y} \in \R^k$ such that $|\mathbf{y}-\mathbf{x}| < r$ (or
$|\mathbf{y}-\mathbf{x}| \le r$).

We call a set $E \subset \R^k$ \emph{convex} if
\[ \lambda\mathbf{x} + (1-\lambda)\mathbf{y} \in E \]
whenever $\mathbf{x} \in E$, $\mathbf{y} \in E$, and $0 < \lambda < 1$.

For example, balls are convex. For if $|\mathbf{y}-\mathbf{x}| < r$,
$|\mathbf{z}-\mathbf{x}| < r$, and $0 < \lambda < 1$, we have
\begin{align*}
  |\lambda\mathbf{y} + (1-\lambda)\mathbf{z} - \mathbf{x}|
    &= |\lambda(\mathbf{y}-\mathbf{x}) + (1-\lambda)(\mathbf{z}-\mathbf{x})|\\
    &\le \lambda|\mathbf{y}-\mathbf{x}| + (1-\lambda)|\mathbf{z}-\mathbf{x}|
       < \lambda r + (1-\lambda)r\\
    &= r.
\end{align*}
The same proof applies to closed balls. It is also easy to see that $k$-cells
are convex.

\begin{deeper}[Convexity is not a metric notion]
Notice that convexity is defined only for $E \subset \R^k$, never for a general
metric space --- and it could not be, because
$\lambda\mathbf{x} + (1-\lambda)\mathbf{y}$ requires addition and scalar
multiplication, which 2.15 does not provide. Convexity belongs to the vector
space structure, not the metric.

Rudin introduces it here and then barely uses it in this chapter --- Exercise
2.21(c) below proves convex sets are connected, and the next real use is
Theorem 9.19, where the mean value inequality needs the segment joining two
points to stay inside the domain. (Convex \emph{functions}, a different
notion, appear in Exercise 4.23.) The line ``balls are convex'' is worth
keeping, though: it depends on the triangle inequality alone and is the reason
pictures of neighbourhoods may be drawn as round blobs without loss.

\medskip
\begin{center}
\begin{tikzpicture}[x=1mm, y=1mm]
  % convex
  \fill[black!8] plot[smooth cycle, tension=0.9]
    coordinates {(0,3) (14,0) (27,4) (25,13) (10,15)};
  \draw[thin] plot[smooth cycle, tension=0.9]
    coordinates {(0,3) (14,0) (27,4) (25,13) (10,15)};
  \node[ptfill, minimum size=2.4pt] (ca) at (5,7) {};
  \node[ptfill, minimum size=2.4pt] (cb) at (22,8) {};
  \draw[thin] (ca) -- (cb);
  \node[lbl, anchor=north] at (13,-3) {convex: every segment stays in};
  % non-convex
  \begin{scope}[xshift=52mm]
  \fill[black!8] plot[smooth cycle, tension=0.9]
    coordinates {(0,1) (28,1) (25,12) (14,5) (3,12)};
  \draw[thin] plot[smooth cycle, tension=0.9]
    coordinates {(0,1) (28,1) (25,12) (14,5) (3,12)};
  \node[ptfill, minimum size=2.4pt] (na) at (4,8) {};
  \node[ptfill, minimum size=2.4pt] (nb) at (24,8) {};
  \draw[guide] (na) -- (nb);
  \node[lbl, anchor=north] at (14,-3) {not convex: this one leaves};
  \end{scope}
\end{tikzpicture}
\end{center}
\figcap{The definition quantifies over every pair of points and every
$\lambda$: one escaping chord is enough to fail it.}
\end{deeper}

\ritem{2.18}{Definition}
Let $X$ be a metric space. All points and sets mentioned below are understood
to be elements and subsets of $X$.
\begin{enumerate}[label=(\alph*), leftmargin=2.2em, itemsep=3pt, topsep=4pt]
  \item A \emph{neighborhood} of $p$ is a set $N_r(p)$ consisting of all $q$
        such that $d(p,q) < r$ for some $r > 0$. The number $r$ is called the
        \emph{radius} of $N_r(p)$.
  \item A point $p$ is a \emph{limit point} of the set $E$ if every
        neighborhood of $p$ contains a point $q \ne p$ such that
        $q \in E$.\kw{The clause $q \ne p$ is the whole definition. Without it
        every point of $E$ would trivially be a limit point of $E$, and the
        notion would be useless. A limit point is a point that $E$ crowds
        \emph{from outside itself} --- and it need not belong to $E$: see
        2.21(e), where $0$ is a limit point of $\{1/n\}$ and lies nowhere in
        the set.}
  \item If $p \in E$ and $p$ is not a limit point of $E$, then $p$ is called an
        \emph{isolated point} of $E$.
  \item $E$ is \emph{closed} if every limit point of $E$ is a point of $E$.
  \item A point $p$ is an \emph{interior point} of $E$ if there is a
        neighborhood $N$ of $p$ such that $N \subset E$.
  \item $E$ is \emph{open} if every point of $E$ is an interior point of $E$.
  \item The \emph{complement} of $E$ (denoted by $E^c$) is the set of all
        points $p \in X$ such that $p \notin E$.
  \item $E$ is \emph{perfect} if $E$ is closed and if every point of $E$ is a
        limit point of $E$.
  \item $E$ is \emph{bounded} if there is a real number $M$ and a point
        $q \in X$ such that $d(p,q) < M$ for all $p \in E$.
  \item $E$ is \emph{dense} in $X$ if every point of $X$ is a limit point of
        $E$, or a point of $E$ (or both).
\end{enumerate}
Let us note that in $\R^1$ neighborhoods are segments, whereas in $\R^2$
neighborhoods are interiors of circles.

\begin{unpack}[The ten definitions in one picture]
Take $E$ to be a shaded blob in the plane together with one stray point off to
the side. Every term in 2.18 can be read off it.

\medskip
\begin{center}
\begin{tikzpicture}[x=1mm, y=1mm]
  % the blob
  \fill[black!10] plot[smooth cycle, tension=0.8]
    coordinates {(6,8) (26,4) (44,12) (40,28) (18,30)};
  \draw[thin] plot[smooth cycle, tension=0.8]
    coordinates {(6,8) (26,4) (44,12) (40,28) (18,30)};
  \node[lbl] at (25,19) {$E$};
  % interior point
  \node[ptfill, minimum size=2.4pt] at (15,14) {};
  \draw[thin] (15,14) circle (4.5mm);
  \draw[thin] (15,9) -- (13,-4);
  \node[lbl, anchor=north] at (13,-5) {interior point:};
  \node[lbl, anchor=north] at (13,-9) {a ball fits inside $E$};
  % limit point on the edge, not in E
  \node[ptopen] at (43,17) {};
  \draw[guide] (43,17) circle (4mm);
  \draw[thin] (43,21) -- (46,33);
  \node[lbl, anchor=south west] at (44,33) {limit point, not in $E$:};
  \node[lbl, anchor=north west] at (44,33) {every ball meets $E$};
  % isolated point, in E
  \node[ptfill, minimum size=2.6pt] at (74,14) {};
  \draw[thin] (74,14) circle (4.5mm);
  \draw[thin] (74,9) -- (74,-4);
  \node[lbl, anchor=north] at (74,-5) {isolated point of $E$:};
  \node[lbl, anchor=north] at (74,-9) {its ball meets $E$};
  \node[lbl, anchor=north] at (74,-13) {only at itself};
  \node[lbl, anchor=south] at (74,20) {also part of $E$};
\end{tikzpicture}
\end{center}
\figcap{A neighbourhood is a ball. An interior point has one inside $E$; a
limit point has every one of its balls meeting $E$ elsewhere; an isolated
point of $E$ has a ball meeting $E$ only at itself.}

\smallskip
Three consequences are worth fixing now, because Rudin uses them without
comment.

\smallskip
\textbf{Open and closed are not opposites.} They are defined by unrelated
conditions --- (f) is about points \emph{of} $E$, (d) is about limit points,
which may lie outside. A set can be both (the whole space $X$, and
$\emptyset$), or neither (the half-open $[0,1)$ in $\R$). Theorem 2.23 supplies
the only link there is: $E$ open $\iff$ $E^c$ closed.

\smallskip
\textbf{A point of $E$ is either isolated or a limit point of $E$}, by (c), and
never both. A perfect set (h) is one with no isolated points at all, closed.

\smallskip
\textbf{``Dense'' means the closure is everything.} Once 2.26 defines
$\overline{E} = E \cup E'$, clause (j) reads $\overline{E} = X$. Theorem
1.20(b) says exactly that $\Q$ is dense in $\R$.
\end{unpack}

\ritem{2.19}{Theorem}
\emph{Every neighborhood is an open set.}

\begin{proof}
Consider a neighborhood $E = N_r(p)$, and let $q$ be any point of $E$. Then
there is a positive real number $h$ such that
\[ d(p,q) = r - h. \]
For all points $s$ such that $d(q,s) < h$, we have then
\[ d(p,s) \le d(p,q) + d(q,s) < r - h + h = r, \]
so that $s \in E$. Thus $q$ is an interior point of $E$.
\end{proof}

\begin{unpack}[A ball inside a ball]
The statement sounds like a tautology and is not: ``neighborhood'' is defined
by a distance inequality, ``open'' by the existence of small balls around every
point. The theorem says the first implies the second, and the triangle
inequality is what does it.

\medskip
\begin{center}
\begin{tikzpicture}[x=1mm, y=1mm]
  \draw[thin] (0,0) circle (18mm);
  \node[ptfill, minimum size=2.4pt] at (0,0) {};
  \node[mlbl, anchor=north east] at (-0.5,-0.5) {$p$};
  \node[ptfill, minimum size=2.4pt] at (11,7) {};
  \node[mlbl, anchor=south] at (11,8.5) {$q$};
  \draw[thin] (11,7) circle (5mm);
  \draw[thin, -{Stealth[length=3pt]}] (0,0) -- (11,7);
  \node[mlbl, anchor=north west] at (4,4) {$r-h$};
  \draw[thin, -{Stealth[length=3pt]}] (11,7) -- (14.5,10.6);
  \node[mlbl, anchor=west] at (16,11) {$h$};
  \node[lbl, anchor=west] at (24,0) {radius $r$ from $p$};
  \node[lbl, anchor=west] at (24,-4) {leaves room $h$ at $q$};
\end{tikzpicture}
\end{center}
\figcap{Whatever is left over after travelling $d(p,q)$ is exactly the radius
that fits at $q$.}

\smallskip
The number $h = r - d(p,q)$ is positive precisely because $q$ is \emph{inside}
the neighbourhood. That is the only thing being used, and it is why closed
balls in $\R^k$ are not open: a point on the boundary sphere has $h = 0$ and
no room at all.
\end{unpack}

\ritem{2.20}{Theorem}
\emph{If $p$ is a limit point of a set $E$, then every neighborhood of $p$
contains infinitely many points of $E$.}

\begin{proof}
Suppose there is a neighborhood $N$ of $p$ which contains only a finite number
of points of $E$. Let $q_1,\dots,q_n$ be those points of $N \cap E$ which are
distinct from $p$, and put
\[ r = \min_{1 \le m \le n} d(p,q_m). \]
[We use this notation to denote the smallest of the numbers
$d(p,q_1)$, \dots, $d(p,q_n)$.] The minimum of a finite set of positive numbers is
clearly positive, so that $r > 0$.\kn{Here is where finiteness is spent, and it
is the entire proof. An infinite set of positive numbers can have infimum $0$
--- take $\{1/n\}$ --- so no such $r$ could be extracted. The same
finite-minimum move appears in 2.24(c), and both are the reason those theorems
say ``finite.''}

The neighborhood $N_r(p)$ contains no point $q$ of $E$ such that $q \ne p$, so
that $p$ is not a limit point of $E$. This contradiction establishes the
theorem.
\end{proof}

\ritem{}{Corollary}
\emph{A finite point set has no limit points.}

\ritem{2.21}{Examples}
Let us consider the following subsets of $\R^2$:
\begin{enumerate}[label=(\alph*), leftmargin=2.2em, itemsep=2pt, topsep=4pt]
  \item The set of all complex $z$ such that $|z| < 1$.
  \item The set of all complex $z$ such that $|z| \le 1$.
  \item A nonempty finite set.
  \item The set of all integers.
  \item The set consisting of the numbers $1/n$ $(n = 1,2,3,\dots)$. Let us
        note that this set $E$ has a limit point (namely, $z = 0$) but that no
        point of $E$ is a limit point of $E$; we wish to stress the difference
        between having a limit point and containing one.
  \item The set of all complex numbers (that is, $\R^2$).
  \item The segment $(a,b)$.
\end{enumerate}
Let us note that (d), (e), (g) can be regarded also as subsets of $\R^1$.

Some properties of these sets are tabulated below:

\medskip
\begin{center}
\begin{tabular}{@{}c@{\hspace{2em}}c@{\hspace{1.6em}}c@{\hspace{1.6em}}c@{\hspace{1.6em}}c@{}}
\toprule
 & Closed & Open & Perfect & Bounded\\
\midrule
(a) & No  & Yes & No  & Yes\\
(b) & Yes & No  & Yes & Yes\\
(c) & Yes & No  & No  & Yes\\
(d) & Yes & No  & No  & No\\
(e) & No  & No  & No  & Yes\\
(f) & Yes & Yes & Yes & No\\
(g) & No  &     & No  & Yes\\
\bottomrule
\end{tabular}
\end{center}
\medskip

In (g), we left the second entry blank. The reason is that the segment $(a,b)$
is not open if we regard it as a subset of $\R^2$, but it is an open subset of
$\R^1$.\kw{The blank entry is the most instructive cell in the table. ``Open''
is not a property of a set; it is a property of a set \emph{in an ambient
space}. As a subset of the line, $(a,b)$ is open; the same points sitting in
the plane form a segment with empty interior, since no disc fits inside it.
This is what 2.29--2.30 are about, and forgetting it is the standard error in
Chapter 4.}

\begin{unpack}[Reading the table]
Row (f) is worth staring at: $\R^2$ is closed \emph{and} open in itself. So is
$\emptyset$. That is not a paradox --- it is what the definitions say. In fact
connectedness (2.45) amounts to exactly this: a space is connected precisely
when $\emptyset$ and the whole space are its only subsets that are both open
and closed, as a short argument from Definition 2.45 shows.

Row (e) is the one Rudin flags. $E = \{1/n\}$ \emph{has} the limit point $0$
but does not \emph{contain} it, so $E$ is not closed. No point of $E$ is a
limit point of $E$, so every point of $E$ is isolated. A set can therefore be
entirely isolated points and still have a limit point elsewhere.

Rows (b) and (f) are the perfect ones. Perfect means closed with no isolated
points, and (c),(d),(e) fail it for having isolated points, while (a) and (g)
fail it for not being closed.
\end{unpack}

\ritem{2.22}{Theorem}
\emph{Let $\{E_\alpha\}$ be a (finite or infinite) collection of sets
$E_\alpha$. Then}
\begin{equation}
  \Bigl(\bigcup_\alpha E_\alpha\Bigr)^c = \bigcap_\alpha (E_\alpha^c). \tag{2.20}
\end{equation}

\begin{proof}
Let $A$ and $B$ be the left and right members of (2.20). If $x \in A$, then
$x \notin \bigcup_\alpha E_\alpha$, hence $x \notin E_\alpha$ for any $\alpha$,
hence $x \in E_\alpha^c$ for every $\alpha$, so that $x \in \bigcap
E_\alpha^c$. Thus $A \subset B$.

Conversely, if $x \in B$, then $x \in E_\alpha^c$ for every $\alpha$, hence
$x \notin E_\alpha$ for any $\alpha$, hence $x \notin \bigcup_\alpha E_\alpha$,
so that $x \in (\bigcup_\alpha E_\alpha)^c$. Thus $B \subset A$.

It follows that $A = B$.
\end{proof}

\ritem{2.23}{Theorem}
\emph{A set $E$ is open if and only if its complement is closed.}

\begin{proof}
First, suppose $E^c$ is closed. Choose $x \in E$. Then $x \notin E^c$, and $x$
is not a limit point of $E^c$. Hence there exists a neighborhood $N$ of $x$
such that $E^c \cap N$ is empty, that is, $N \subset E$. Thus $x$ is an
interior point of $E$, and $E$ is open.

Next, suppose $E$ is open. Let $x$ be a limit point of $E^c$. Then every
neighborhood of $x$ contains a point of $E^c$, so that $x$ is not an interior
point of $E$. Since $E$ is open, this means that $x \in E^c$. It follows that
$E^c$ is closed.
\end{proof}

\ritem{}{Corollary}
\emph{A set $F$ is closed if and only if its complement is open.}

\begin{deeper}[2.22 and 2.23 are the engine of the section]
Every remaining proof about closed sets is obtained by complementing a proof
about open sets. 2.22 (De Morgan) turns unions into intersections; 2.23 turns
open into closed. Together they let Rudin prove half of 2.24 and get the other
half free, which is exactly what he does.

The pattern is worth naming because it recurs throughout the book: prove the
statement for open sets, where the definition is local and easy to verify at a
point, then complement. Closedness is a statement about limit points, which are
harder to handle directly --- so one almost never argues with it head on.
\end{deeper}

\ritem{2.24}{Theorem}
\begin{enumerate}[label=(\alph*), leftmargin=2.2em, itemsep=1pt, topsep=3pt]
  \item \emph{For any collection $\{G_\alpha\}$ of open sets,
        $\bigcup_\alpha G_\alpha$ is open.}
  \item \emph{For any collection $\{F_\alpha\}$ of closed sets,
        $\bigcap_\alpha F_\alpha$ is closed.}
  \item \emph{For any finite collection $G_1,\dots,G_n$ of open sets,
        $\bigcap_{i=1}^n G_i$ is open.}
  \item \emph{For any finite collection $F_1,\dots,F_n$ of closed sets,
        $\bigcup_{i=1}^n F_i$ is closed.}
\end{enumerate}

\begin{proof}
Put $G = \bigcup_\alpha G_\alpha$. If $x \in G$ then $x \in G_\alpha$ for some
$\alpha$. Since $x$ is an interior point of $G_\alpha$, $x$ is also an interior
point of $G$, and $G$ is open. This proves (a).

By Theorem 2.22,
\begin{equation}
  \Bigl(\bigcap_\alpha F_\alpha\Bigr)^c = \bigcup_\alpha (F_\alpha^c),
  \tag{2.21}
\end{equation}
and $F_\alpha^c$ is open, by Theorem 2.23. Hence (a) implies that (2.21) is
open, so that $\bigcap_\alpha F_\alpha$ is closed.

Next, put $H = \bigcap_{i=1}^n G_i$. For any $x \in H$, there exist
neighborhoods $N_i$ of $x$, with radii $r_i$, such that $N_i \subset G_i$
$(i = 1,\dots,n)$. Put
\[ r = \min\{r_1,\dots,r_n\} \]
and let $N$ be the neighborhood of $x$ of radius $r$. Then $N \subset G_i$ for
$i = 1,\dots,n$, so that $N \subset H$, and $H$ is open.\kn{The finite minimum
again, exactly as in 2.20. With infinitely many $G_i$ the radii could shrink to
zero and leave no room, which is precisely what happens in Example 2.25.}

By taking complements, (d) follows from (c):
\[ \Bigl(\bigcup_{i=1}^n F_i\Bigr)^c = \bigcap_{i=1}^n (F_i^c). \qedhere \]
\end{proof}

\ritem{2.25}{Example}
In parts (c) and (d) of the preceding theorem, the finiteness of the
collections is essential. For let $G_n$ be the segment $(-1/n,\,1/n)$
$(n = 1,2,3,\dots)$. Then $G_n$ is an open subset of $\R^1$. Put
$G = \bigcap_{n=1}^\infty G_n$. Then $G$ consists of a single point (namely,
$x = 0$) and is therefore not an open subset of $\R^1$.

Thus the intersection of an infinite collection of open sets need not be open.
Similarly, the union of an infinite collection of closed sets need not be
closed.

\begin{center}
\begin{tikzpicture}[x=1mm, y=1mm]
  \foreach \n/\w in {1/40, 2/20, 3/13.3, 4/10}{
    \pgfmathsetmacro\yy{-(\n-1)*7}
    \draw[axis] (-46,\yy) -- (46,\yy);
    \draw[seg] (-\w,\yy) -- (\w,\yy);
    \node[ptopen] at (-\w,\yy) {}; \node[ptopen] at (\w,\yy) {};
    \node[lbl, anchor=west] at (50,\yy) {$G_{\n}$};}
  \draw[axis] (-46,-33) -- (46,-33);
  \node[ptfill, minimum size=3pt] at (0,-33) {};
  \node[lbl, anchor=west] at (50,-33) {$\bigcap G_n=\{0\}$};
  \draw[guide] (0,3) -- (0,-30);
\end{tikzpicture}
\end{center}
\figcap{Each $G_n$ is open, yet the radii shrink to nothing and the
intersection is a single point --- not open.}

\ritem{2.26}{Definition}
If $X$ is a metric space, if $E \subset X$, and if $E'$ denotes the set of all
limit points of $E$ in $X$, then the \emph{closure} of $E$ is the set
$\overline{E} = E \cup E'$.

\ritem{2.27}{Theorem}
\emph{If $X$ is a metric space and $E \subset X$, then}
\begin{enumerate}[label=(\alph*), leftmargin=2.2em, itemsep=1pt, topsep=3pt]
  \item \emph{$\overline{E}$ is closed,}
  \item \emph{$E = \overline{E}$ if and only if $E$ is closed,}
  \item \emph{$\overline{E} \subset F$ for every closed set $F \subset X$ such
        that $E \subset F$.}
\end{enumerate}
\emph{By (a) and (c), $\overline{E}$ is the smallest closed subset of $X$ that
contains $E$.}\kn{That last sentence is the content; (a) and (b) are warm-up.
``Smallest closed set containing $E$'' is the characterisation you should carry
--- it is how closure is \emph{defined} in most later treatments, with
$\overline E = E \cup E'$ demoted to a theorem.}

\begin{proof}
\begin{enumerate}[label=(\alph*), leftmargin=2.2em, itemsep=4pt, topsep=3pt]
  \item If $p \in X$ and $p \notin \overline{E}$ then $p$ is neither a point of
        $E$ nor a limit point of $E$. Hence $p$ has a neighborhood which does
        not intersect $E$. The complement of $\overline{E}$ is therefore open.
        Hence $\overline{E}$ is closed.
  \item If $E = \overline{E}$, (a) implies that $E$ is closed. If $E$ is
        closed, then $E' \subset E$ [by Definitions 2.18(d) and 2.26], hence
        $E = \overline{E}$.
  \item If $F$ is closed and $F \supset E$, then $F \supset F'$, hence
        $F \supset E'$. Thus $F \supset \overline{E}$.
\end{enumerate}
\end{proof}

\begin{center}
\begin{tikzpicture}[x=1mm, y=1mm]
  % E = {1/n}
  \draw[axis] (-2,10) -- (72,10);
  \foreach \n in {1,...,9}{\node[ptfill, minimum size=2pt] at ({62/\n},10) {};}
  \node[ptgap] at (0,10) {};
  \draw[guide] (0,12.5) -- (0,17);
  \node[lbl, anchor=south west] at (-2,17) {closure adds the limit point $0$};
  \node[mlbl, anchor=north] at (62,8) {$1$};
  \node[lbl, anchor=east] at (-6,10) {$E=\{1/n\}$};
  % E = (0,1)
  \draw[axis] (-2,-8) -- (72,-8);
  \draw[seg] (12,-8) -- (57,-8);
  \node[ptgap] at (12,-8) {};
  \node[ptgap] at (57,-8) {};
  \draw[guide] (12,-10.5) -- (12,-15);
  \draw[guide] (57,-10.5) -- (57,-15);
  \node[lbl, anchor=north] at (34,-15.5) {closure adds both endpoints};
  \node[lbl, anchor=east] at (-6,-8) {$E=(0,1)$};
\end{tikzpicture}
\end{center}
\figcap{Closure fills in exactly the limit points (dashed circles); nothing
else changes. Compare 2.28: for a bounded set of reals, $\sup E$ is always
among the points filled in, if it was missing.}

\ritem{2.28}{Theorem}
\emph{Let $E$ be a nonempty set of real numbers which is bounded above. Let
$y = \sup E$. Then $y \in \overline{E}$. Hence $y \in E$ if $E$ is closed.}

Compare this with the examples in Sec.\ 1.9.

\begin{proof}
If $y \in E$ then $y \in \overline{E}$. Assume $y \notin E$. For every $h > 0$
there exists then a point $x \in E$ such that $y - h < x < y$, for otherwise
$y - h$ would be an upper bound of $E$. Thus $y$ is a limit point of $E$. Hence
$y \in \overline{E}$.
\end{proof}

\begin{unpack}[Chapter 1 meets Chapter 2]
This is the first theorem to connect the order structure of $\R$ with the
topology, and the proof is two lines because the work was done in Chapter 1.
``For every $h>0$ there is $x \in E$ with $y - h < x$'' is precisely clause
(ii) of Definition 1.8, in the working form given there --- and clause (ii)
says exactly that $y$ is crowded from below by $E$, which is what a limit point
is.

So $\sup E$ is either a member of $E$ or a limit point of it, never anything
else. The examples at 1.9 now read topologically: $\sup E_1 = 0 \notin E_1$
makes $0$ a limit point, while $\sup E_2 = 0 \in E_2$. Theorem 2.28 is what
lets Chapter 4 conclude that a continuous function on a compact set attains its
supremum (4.16).
\end{unpack}

\ritem{2.29}{Remark}
Suppose $E \subset Y \subset X$, where $X$ is a metric space. To say that $E$
is an open subset of $X$ means that to each point $p \in E$ there is associated
a positive number $r$ such that the conditions $d(p,q) < r$, $q \in X$ imply
that $q \in E$. But we have already observed (Example 2.16) that $Y$ is also a
metric space, so that our definitions may equally well be made within $Y$. To
be quite explicit, let us say that $E$ is \emph{open relative to $Y$} if to
each $p \in E$ there is associated an $r > 0$ such that $q \in E$ whenever
$d(p,q) < r$ and $q \in Y$. Example 2.21(g) showed that a set may be open
relative to $Y$ without being an open subset of $X$. However, there is a simple
relation between these concepts, which we now state.

\ritem{2.30}{Theorem}
\emph{Suppose $Y \subset X$. A subset $E$ of $Y$ is open relative to $Y$ if and
only if $E = Y \cap G$ for some open subset $G$ of $X$.}

\begin{proof}
Suppose $E$ is open relative to $Y$. To each $p \in E$ there is a positive
number $r_p$ such that the conditions $d(p,q) < r_p$, $q \in Y$ imply that
$q \in E$. Let $V_p$ be the set of all $q \in X$ such that $d(p,q) < r_p$, and
define
\[ G = \bigcup_{p \in E} V_p. \]
Then $G$ is an open subset of $X$, by Theorems 2.19 and 2.24.

Since $p \in V_p$ for all $p \in E$, it is clear that $E \subset G \cap Y$.

By our choice of $V_p$, we have $V_p \cap Y \subset E$ for every $p \in E$, so
that $G \cap Y \subset E$. Thus $E = G \cap Y$, and one half of the theorem is
proved.

Conversely, if $G$ is open in $X$ and $E = G \cap Y$, every $p \in E$ has a
neighborhood $V_p \subset G$. Then $V_p \cap Y \subset E$, so that $E$ is open
relative to $Y$.
\end{proof}

\begin{center}
\begin{tikzpicture}[x=1mm, y=1mm]
  % ambient X
  \draw[thin] (0,0) rectangle (56,30);
  \node[lbl, anchor=north west] at (1,29.5) {$X=\R^2$};
  % Y = a line inside X
  \draw[seg] (4,12) -- (52,12);
  \node[lbl, anchor=west] at (58,12) {$Y=\R^1$};
  % G an open disc in X
  \draw[guide] (28,15) circle (11mm);
  \node[lbl, anchor=south] at (28,27) {$G$ open in $X$};
  % E = G cap Y
  \draw[draw=RuInk, line width=2pt] (18.5,12) -- (37.5,12);
  \node[lbl, anchor=north] at (28,9) {$E=G\cap Y$: open in $Y$, not in $X$};
\end{tikzpicture}
\end{center}
\figcap{The same points form an open set inside the line and a non-open set
inside the plane. ``Open'' always means \emph{open in} something.}
